% ============================================================
%  CHAPTER 4 — Simulation Framework
%  Logical chain: Proposed Solution → Implementation
%  Status: Not started | Blocking: PQ1 (robot model), PQ2 (simulator)
% ============================================================

\chapter{Simulation Framework}
\label{chap:simulation}


% TODO: Update once PQ1 and PQ2 are resolved (master_memory §16).
% Current candidates: MuJoCo or Isaac Sim; TALOS or custom URDF humanoid.

\lipsum[25]

% ─── 4.1 SIMULATOR SELECTION ────────────────────────────────
\section{Simulator Selection and Justification}


% Compare MuJoCo vs. Isaac Sim:
%   - MuJoCo: fast, stable contacts (soft model), widely used in WBC literature
%   - Isaac Sim: GPU-accelerated, photorealistic, better for parallel training
% Decision criteria: contact fidelity, reproducibility, community support.
% TODO: Fill in final choice + rationale once PQ2 resolved.

\lipsum[26]

\begin{criticalinsight}
  % TODO: Justify how the simulator choice affects the validity of conclusions
  % (§6 Problem 1, §8.4 Risk row 4).
\end{criticalinsight}

% ─── 4.2 ROBOT MODEL ─────────────────────────────────────────
\section{Robot Model}


% TODO: Fill in once PQ1 resolved (TALOS / Unitree H1 / custom URDF).
% Document: DOF count, actuation type (SEA / ideal torque),
% sensor suite (IMU, joint encoders, foot F/T), URDF source.

\begin{table}[htb]
  \centering
  
  
  \begin{tabularx}{\linewidth}{lX}
    \toprule
    \textbf{Property}   & \textbf{Value} \\
    \midrule
    Platform            & [TALOS / Unitree H1 / Custom] \\
    Total DOF           & [28–32] \\
    Actuation           & [Torque-controlled / SEA] \\
    Base sensors        & IMU (1,kHz) \\
    Joint sensors       & Encoders (1,kHz) \\
    Foot sensors        & Force/torque (1,kHz) \\
    URDF source         & [URL / repository] \\
    \bottomrule
  \end{tabularx}
\end{table}

\lipsum[27]

\begin{criticalinsight}
  % TODO: Discuss URDF fidelity as a proxy for hardware accuracy
  % (actuator model, damping, friction parameters).
\end{criticalinsight}

% ─── 4.3 SOFTWARE STACK ─────────────────────────────────────
\section{Software Architecture}


% Components (master_memory §8.2):
%   - Pinocchio: rigid-body dynamics (RNEA, Jacobians, ABA)
%   - OSQP / qpOASES: QP solver (warm-startable, real-time capable)
%   - Crocoddyl: centroidal MPC (optional, tie to novelty question PQ4)
%   - Simulator interface: Python bindings or C++ ROS2 node
%   - ROS2: optional middleware layer

% TODO: Insert software architecture diagram (figures/software_stack.pdf)

\begin{figure}[htb]
  \centering
  \fbox{\rule{0pt}{4cm}\rule{8cm}{0pt}}  % placeholder
  \caption{Software stack architecture. Real-time loop runs at 1,kHz;
           MPC runs asynchronously at 1,kHz.}
  
\end{figure}

\lipsum[28]

\begin{criticalinsight}
  % TODO: Discuss real-time constraints: QP solve time < 1 ms per cycle
  % (master_memory §2.3). Benchmark solver against constraint.
\end{criticalinsight}

% ─── 4.4 CONTROL LOOP IMPLEMENTATION ────────────────────────
\section{Control Loop Implementation}


% Describe:
%   - Sensor data acquisition → state estimation (EKF)
%   - MPC update (asynchronous thread)
%   - WBC QP formulation → solve → torque dispatch
%   - Contact mode detection (threshold / probabilistic)

\begin{verbatim}[H]
  
  
  
    Input: Sensor readings $(\q, \dq, \fc^{\mathrm{meas}})$, MPC reference $\xtask^{\mathrm{ref}}$
    Output: Joint torques $\torque$
     Estimate floating-base state via EKF: $\hat{\q}_b, \dot{\hat{\q}}_b$
     Detect contact modes from foot force readings
     Update contact Jacobian $\Jc$ and task Jacobians $\J_i$
     Formulate QP \cref{eq:qp-cost}--\cref{eq:qp-friction}
     Solve QP $\rightarrow$ $(\ddq^*, \torque^*, \fc^*)$
     return $\torque^*$
  
\end{verbatim}

\lipsum[29]

% ─── 4.5 VALIDATION OF SIMULATION SETUP ─────────────────────
\section{Validation of the Simulation Setup}


% Validate against published benchmark scenarios before running S1–S5.
% (master_memory §8.4 Risk row 4 mitigation).
% e.g., reproduce a standing balance result from Koolen et al. (2016).

\lipsum[30]

\begin{criticalinsight}
  % TODO: Identify which published results serve as ground truth for the
  % simulation validation step, and justify the choice.
\end{criticalinsight}
